% part: intuitionistic-logic
% chapter: soundness-completeness
% section: truth-lemma

\documentclass[../../../include/open-logic-section]{subfiles}

\begin{document}

\olfileid{int}{sc}{tru}

\olsection{لمّة الصدق}

هذه اللمّة تبين الصلة بين الإشباع في النموذج القانوني وبين كون
ال!!{formula}s من !!{element}s المجموعة الأولية~$\Delta$.

\begin{lem}\ollabel{lem:truth}
  متى كانت~$\Delta$ أولية، تحقق~$\mSat{M(\Delta)}{!A}[\sigma]$ إذا
  وفقط إذا تحقق~$\Delta(\sigma) \Proves !A$.
\end{lem}

\begin{proof}
  نستقرئ في تركيب~$!A$.
  \begin{enumerate}
    \item \indcase{!A}{\lfalse}{أولية~$\Delta(\sigma)$ تقتضي
      اتساقها؛ فيثبت~$\Delta(\sigma) \Proves/ \indfrm$. والتعريف يقضي بـ
      $\mSat/{M(\Delta)}{\indfrm}[\sigma]$.}
    \item \indcase{!A}{p}{من تعريف~$\mSat{{}}{}$ نعلم أن
      $\mSat{M(\Delta)}{\indfrm}[\sigma]$ إذا وفقط إذا كان
      $\sigma \in V(p)$، وهذا يكافئ
      $\Delta(\sigma) \Proves \indfrm$.}
    \item \indcase{!A}{\lnot !B}{وجه هذه الحالة كوجه الشرط الذي نتيجته
      $\lfalse$، وتفصيله فيما يأتي. ليكن أولًا
      $\Delta(\sigma)\Proves\lnot !B$. فكل~$\sigma'$ يحقق
      $R\sigma\sigma'$ يحقق كذلك
      $\Delta(\sigma)\subseteq\Delta(\sigma')$، فيلزم
      $\Delta(\sigma')\Proves\lnot !B$. ولو تحقق
      $\mSat{M(\Delta)}{!B}[\sigma']$، لثبت بالاستقراء
      $\Delta(\sigma')\Proves !B$، وذلك يناقض اتساق
      $\Delta(\sigma')$. فلا يصدق~$!B$ في امتداد~$\sigma'$ أصلًا،
      ومعناه~$\mSat{M(\Delta)}{\lnot !B}[\sigma]$.

      وأما عكس هذا الاتجاه فنقيمه بطريق النقيض المقابل. ليكن
      $\Delta(\sigma)\Proves/\lnot !B$؛ فيلزم
      $\Delta(\sigma)\cup\{!B\}\Proves/\lfalse$. إذ لو اشتُق
      $\lfalse$ بعد ضم~$!B$، لأعطى إدخال النفي
      $\lnot !B$ من~$\Delta(\sigma)$. ونأخذ
      $\tuple{!B,\lfalse}=\tuple{!B_n,!C_n}$ لبعض~$n$. فيكون
      $\Delta(\sigma.n)=(\Delta(\sigma)\cup\{!B\})^*$، ويلزم
      $\Delta(\sigma.n)\Proves !B$. فبالاستقراء
      $\mSat{M(\Delta)}{!B}[\sigma.n]$؛ ومع
      $R\sigma(\sigma.n)$ يثبت
      $\mSat/{M(\Delta)}{\lnot !B}[\sigma]$.}
    \item \indcase{!A}{!B \land
      !C}{$\mSat{M(\Delta)}{\indfrm}[\sigma]$ يكافئ اجتماع
      $\mSat{M(\Delta)}{!B}[\sigma]$ و
      $\mSat{M(\Delta)}{!C}[\sigma]$. وبالاستقراء
      $\mSat{M(\Delta)}{!B}[\sigma]$ إذا وفقط إذا كان
      $\Delta(\sigma) \Proves !B$، ومثله في~$!C$. واجتماع
      $\Delta(\sigma) \Proves !B$ و$\Delta(\sigma) \Proves !C$ يكافئ
      $\Delta(\sigma) \Proves \indfrm$.}
    \item \indcase{!A}{!B \lor
      !C}{$\mSat{M(\Delta)}{\indfrm}[\sigma]$ يكافئ أحد الأمرين:
      $\mSat{M(\Delta)}{!B}[\sigma]$ أو
      $\mSat{M(\Delta)}{!C}[\sigma]$. وهما بالاستقراء يكافئان
      $\Delta(\sigma) \Proves !B$ أو
      $\Delta(\sigma) \Proves !C$. فلا يبقى إلا أن نبين أن هذين يكافئان
      $\Delta(\sigma) \Proves \indfrm$. أما الذهاب من اليسار إلى
      اليمين فبيّن، وأما الرجوع من اليمين إلى اليسار فمبناه
      أولية~$\Delta(\sigma)$.}
    \item \indcase{!A}{!B \lif !C}{نبدأ بإثبات النقيض المقابل للاتجاه
      من اليسار إلى اليمين. ليكن
      $\Delta(\sigma) \Proves/ !B \lif !C$؛ فيلزم
      $\Delta(\sigma) \cup \{!B\} \Proves/ !C$. والزوج
      $\tuple{!B, !C}$ هو~$\tuple{!B_n, !C_n}$ لبعض~$n$، فيكون
      $\Delta(\sigma.n) = (\Delta(\sigma) \cup \{!B\})^*$، ويثبت
      $\Delta(\sigma.n) \Proves !B$ مع
      $\Delta(\sigma.n) \Proves/ !C$. وبالاستقراء نجد
      $\mSat{M(\Delta)}{!B}[\sigma.n]$ و
      $\mSat/{M(\Delta)}{!C}[\sigma.n]$. ولما كان
      $R\sigma(\sigma.n)$، لزم
      $\mSat/{M(\Delta)}{\indfrm}[\sigma]$.

      ولنأخذ الآن~$\Delta(\sigma) \Proves !B \lif !C$، مع
      $R\sigma\sigma'$. فبما أن
      $\Delta(\sigma) \subseteq \Delta(\sigma')$، يلزم من
      $\Delta(\sigma') \Proves !B$ أن
      $\Delta(\sigma') \Proves !C$. أي إن كل~$\sigma'$ يحقق
      $R\sigma\sigma'$ يحقق أحد الأمرين:~$\Delta(\sigma') \Proves/ !B$ أو
      $\Delta(\sigma') \Proves !C$. وبفرض الاستقراء، متى كان
      $R\sigma\sigma'$ كان إما
      $\mSat/{M(\Delta)}{!B}[\sigma']$ وإما
      $\mSat{M(\Delta)}{!C}[\sigma']$؛ فثبت
      $\mSat{M(\Delta)}{\indfrm}[\sigma]$.}
  \end{enumerate}
\end{proof}

\end{document}
